\documentclass{article}
\usepackage{arxiv}

\usepackage[utf8]{inputenc}
\usepackage[russian, english]{babel}
\usepackage[T1]{fontenc}
\usepackage{url}
\usepackage{booktabs}
\usepackage{amsfonts}
\usepackage{nicefrac}
\usepackage{microtype}
\usepackage{graphicx}
\usepackage{natbib}
\usepackage{doi}



\title{Epiplexity on one GPU: estimators for small networks and when they agree}

\author{
	Alexey Kravatskiy \\
	Moscow Institute of Physics and Technology \\
	\And
	Mark Ikonnikov \\
	Moscow Institute of Physics and Technology \\
	\AND
	Daniil Kazachkov \\
	Moscow Institute of Physics and Technology \\
	\And
	Danila Chernousov \\
	Moscow Institute of Physics and Technology \\
}
\date{}

\renewcommand{\shorttitle}{Epiplexity on one GPU}

%%% Add PDF metadata to help others organize their library
%%% Once the PDF is generated, you can check the metadata with
%%% $ pdfinfo main.pdf
\hypersetup{
pdftitle={Epiplexity on one GPU: estimators for small networks and when they agree},
pdfsubject={cs.LG},
pdfauthor={Alexey Kravatskiy, Mark Ikonnikov, Daniil Kazachkov, Danila Chernousov},
pdfkeywords={epiplexity, minimum description length, prequential coding},
}

\begin{document}
\maketitle

\begin{abstract}
	Epiplexity~\citep{finzi2026epiplexity} measures the structure that a computationally bounded
	learner can extract from data. We estimate epiplexity with small models (MLPs and
	convolutional networks) on small datasets on a single GPU. The \texttt{epimeter} library
	gives one interface that returns the model part and the data part of a code length in bits
	for the prequential, requential, variational and Laplace codes, the classical criteria, EDL
	and the reservoir score. With it we estimate epiplexity on elementary cellular automata,
	label noise, permuted pixels and a planted teacher--student task, measure the spread over
	seeds under a fixed protocol, and find when Bayesian model parts agree with epiplexity.
\end{abstract}


\keywords{Epiplexity \and Minimum description length \and Prequential coding}

% Section plan follows plan/README.md (goal, research questions RQ1--RQ4, protocol).

\section{Introduction}

\section{Background}

\section{Estimators}
\subsection{Prequential code}
\subsection{Requential code}
\subsection{Bayesian codes}
\subsection{Classical criteria and fast proxies}

\section{Experimental setup}

\section{Results}
\subsection{Elementary cellular automata (RQ1)}
\subsection{Label noise (RQ2)}
\subsection{Bayesian model parts (RQ3)}
\subsection{Seed variance and fast proxies (RQ4)}

\section{Conclusion}


\bibliographystyle{unsrtnat}
\bibliography{references}

\end{document}
